\scanpage{51}
\begin{flushright}
atkin \mspage{11}
\end{flushright}

\[
c(29n)=a\,\handwritten{\lambda}(n)+\mathord{\rlap{$\handwritten{\bar{\phantom{a}}}$}a}\,\handwritten{\bar{\lambda}}(n),
\]
\uline{where} $a,\ \handwritten{\lambda}(n)$,\uline{and} \uline{their} \uline{conjugates} $\mathord{\rlap{$\handwritten{\bar{\phantom{a}}}$}a},\ \handwritten{\bar{\lambda}}(n)$ \uline{lie} \uline{in} $\rlap{\smash{\handwritten{\tikz[x=1em,y=1ex,baseline=0pt]{\draw (.09,-.42) .. controls (.48,-.71) and (.92,-.63) .. (1.1,-.28); \draw (.16,-.94)--(.20,-.67) (.03,1.17)--(.20,1.47); \node[inner sep=0pt] at (.87,1.95){\tiny\unsure{?}}; \draw (1.17,2.20)--(1.30,2.58);}}}}\unsure{\mathbb{GF}|\mathrm{GF}}(29^{2})$,\margin{\handwritten{\tikz[baseline=(f.base)]{\node[inner sep=1.5pt](f){$GF$};\draw (f.center) ellipse [x radius=1.12em,y radius=.78em];}}}
\[
\text{\uline{and} }\handwritten{\lambda}(29n)=\handwritten{\lambda}(29)\handwritten{\lambda}(n)\ ,\
\handwritten{\bar{\lambda}}(29n)=\handwritten{\bar{\lambda}}(29)\handwritten{\bar{\lambda}}(n).
\tag*{(5.2)}
\]

Of course $c(29n)$ itself lies in $\mathrm{GF}(29)$. We may infer from (5.2) that \handwritten{,} for all $n\handwritten{\geqslant}1$\handwritten{,}
\[
c(29^{29^{2}}n)\handwritten{\equiv}c(29n)\qquad(\mathrm{mod}\;29).
\]
We have also an M-result for $\handwritten{\lambda}(n)$ of the usual kind:
\[
\handwritten{\lambda}(np)-\handwritten{\lambda}(n)\handwritten{\lambda}(p)+p^{-1}\handwritten{\lambda}(n/p)=0\quad\text{\uline{if} }p\mathrel{\rlap{$\handwritten{/}$}=}29.
\tag*{(5.3)}
\]

It is now quite \unsure{éasy|easy} to foresee the natural general conjecture for powers of $\unsure{29'|29}$ in the form:

\uline{For} \uline{all} $\alpha\handwritten{\geqslant}1$ \uline{we} \uline{have} \handwritten{\tikz[x=1em,y=1em,baseline=0pt]{\draw (0,.45)--(.14,.49);}}
\[
c(29^{\alpha}n)=a\,\handwritten{\lambda}(n)+\mathord{\rlap{$\handwritten{\bar{\phantom{a}}}$}a}\,\handwritten{\bar{\lambda}}(n),
\]
\uline{where} $a,\handwritten{\lambda}(n)$,\uline{and their conjugates} $\mathord{\rlap{$\handwritten{\bar{\phantom{a}}}$}a},\handwritten{\bar{\lambda}}(n)$,\uline{lie in an} \uline{extension} \uline{of} \uline{the} \uline{residue} \uline{class} \uline{ring} \uline{modulo} $29^{\alpha}$ \uline{defined} \uline{by} \uline{an} \uline{irreducible} \uline{quadratic} \uline{equation},\uline{and} \uline{for} $n\handwritten{\geqslant}1$ \uline{we} \uline{have}
\[
\resizebox{.86\linewidth}{!}{$\displaystyle \begin{aligned}[b]
&\handwritten{\lambda}(29n)=\handwritten{\lambda}(29)\handwritten{\lambda}(n)\\
\text{\uline{and} }\;&\handwritten{\lambda}(np)-\handwritten{\lambda}(n)\handwritten{\lambda}(p)+p^{-1}\handwritten{\lambda}(n/p)=0\quad\text{if }p\mathrel{\rlap{$\handwritten{/}$}=}29.
\end{aligned}$}
\tag*{(5.4)}
\]

This is subject to the same reservations as the cases $q=17,19,\text{and }23$; it is only proved for $\alpha\handwritten{\leqslant}5$,and then only by a combination of modular theory and computation.

The case $q=31$ is precisely analogous to $q=29$,while for $q=37$ we find that \uline{three} conjugates are required over an extension of the ring modulo $37^{\alpha}$ defined by an irreducible \uline{cubic} \tikz[x=1em,y=1em,baseline=0pt]{\fill (.03,.34) circle[radius=.028];} equation.At this point it seemed unprofitable to compute any further cases in detail, since no essentially new features arose after $29$. However brief pre\inserted{\handwritten{l}}iminary investigations for all primes \uline{$q$} with $41\handwritten{\leqslant}q\handwritten{\leqslant}67$ suggest that in general $c(q^{\alpha}n)$ will be a linear combination of $\handwritten{[}q/12\handwritten{]}$ \inserted{\handwritten{multiplicative}} coefficients lying in various extensions of the ring modulo $q^{\alpha}$. In addition the analogy with Hecke theory for modular forms of negative dimension
% archive sheet number: 50
